
\section{Results}
\label{results}
We now demonstrate parameter tuning and propagation metric selection on both Xilinx and Intel systems. We perform our classification experiment on Xilinx systems to determine if the co-tenant is a cryptographic core and, if so, perform a correlation power analysis. The Xilinx sensor classification data on 13 applications, and classifier network are released as open source.

%the pipeline of Figure~\ref{fig:intro_fig}. This includes performing a classification experiment to determine if the co-tenant is a cryptographic core and, if so, performing a Correlation Power Analysis attack to extract a key. The sensor, classification data on 13 applications, and classifier network are all released as open source. 

%on measurements of the \name TDC and their impact on classifying 13 co-located application circuits.
%In this section, we report results on the impact of $\theta$, $\phi$, and propagation metrics on measurements of the \name TDC and their impact on classifying 13 co-located application circuits.

\begin{comment}
Section~\ref{sensor_config} reports the experimental platforms and sensor configuration parameters.
Section~\ref{applications} describes the 13 application circuits that we attempt to classify. 
Section~\ref{capture_clock_tuning} demonstrates the behavior of $\theta$ tuning, including an analysis of delay line properties, how the underlying CARRY primitives must be considered when using a TDC, and how the metrics of propagation presented in Section~\ref{propagation_metric} have different behaviors on CARRY primitive boundaries.
%Section~\ref{channel_capacity} introduces $\phi$ tuning and studies how the three metrics of propagation affect side-channel capacity.
Section~\ref{target_clock_tuning} demonstrates how sweeping $\phi$ can be used to tune the target computation clock and why background subtraction is necessary to tune to the target computation.
Finally, in Section~\ref{classification}, we combine these techniques in a realistic TDC use case where an attacker must extract the architecture and co-tenant computation on a previously unseen FPGA system. \dr{Could cut at the end}
\end{comment}

%\cd{TODO-I need to standardize around delay line rather than carry-chain}
\subsection{Experimental Setup}
\label{sensor_config}

Our experimental platforms for threat model exploration are Amazon Web Services (AWS) EC2 F1 instances with Xilinx UltraScale+ XCVU9P-FLGB2104 FPGAs and six PYNQ-Z2 boards with Xilinx ZYNQ XC7Z020-1CLG400C FPGAs. For an exploration of sensor portability to Intel FPGAs, we target the Terasic DE10-Nano development kit with an Intel Cyclone V SE 5CSEBA6U23I7 FPGA. The ZYNQ XC7Z020 SoC (28nm, 85K LE, dual-core ARM A9 SoC) and the Cyclone V SE 5CSEBA6U23I7 (28nm, 110K LE, dual-core ARM A9 SoC), share similar device specifications and performance allowing for a more direct comparison of delay line features between the Pynq Z2 and DE10-Nano platforms.

On the PYNQ systems, the device is programmed with our sensor and test designs through the Python Productivity for Zynq (PYNQ) infrastructure. The AWS EC2 F1 instances are launched through the EC2 interface and programmed with the unique AGFI identifier associated with our sensor designs. The AGFI is generated by Amazon's unmodified compilation flow with the design checkpoint we provide. Our sensor has passed all design analysis techniques performed by AWS. A 64-bit \name TDC is instantiated on PYNQ-Z2 and a 256-bit \name TDC on AWS. The launch and capture clock domains operate at 100 MHz. This results in a sampling rate, $F_{sample}$, of 25 MHz. MMCM \textcircled{1}, which allows for the phase shifting of $F_{sample}$, produces a 100MHz output clock. The internal $F_{vco}$ is maximized for the two MMCMs so that the step granularity of $\theta$ and $\phi$ is maximized with a step size of 11.16~ps on AWS and 14.88~ps on PYNQ. 

The Intel DE10-Nano system is programmed with a 128-bit instantiation of our sensor, a hard processor system IP that interfaces the Cylonce V FPGA to the co-packaged dual-core ARM A9 hard processor, and a modular scatter-gather DMA engine to offload sensor readings to the system memory. The launch and capture domains operate at 100MHz ($F_{sample}=25Mhz$). Instead of maximizing the internal $F_{vco}$ of the PLLs within the clock generator modules, we improve phase step granularity using the PLL configuration prepossessing algorithm described in Section \ref{sensor_structure}. The output PLL configurations of the algorithm are driven to the Intel PLL reconfiguration core in clock generator \textcircled{1} as consecutive coarse-grained phase sweeps. The measurements from each coarse grained sweep are aggregated into a composite dataset that is sorted with respect to the total phase offset of each measurement. This procedure results in an effective phase-step granularity of 1.6ps as shown in Figure \ref{fig:phase-shift-intel}.

%\dr{Please double check}\cd{1/(56*(1200MHz) = 14.88ps on PYNQ and 1/(56*(1600MHz)) = 11.16 on AWS}

%\subsubsection*{Applications}
\subsection{Applications}
\label{applications}

Our experiments on Intel systems showcase the portability of our sensor and demonstrate the capability of our novel  PLL configuration prepossessing algorithm. on Xilinx systems we extend the analysis of our \name TDC to classify the characteristics of a co-tenant. We have 13 unique applications containing IP cores using different architectural features. The application IP core and the sensor are implemented on the same FPGA. They are logically and physically isolated. The characteristics of the applications are described in the following paragraphs.
%Each of the secrets is instantiated within the same FPGA fabric but logically isolated from our \name TDC. We deploy 13 unique combinations as described below: 

%\paragraph{Sensor Only:} 
\textbf{Intel and Xilinx Sensor Only: } The primary goal of the sensor-only design is to model the lack of another co-tenant. This design only contains the \sensor and associated data collection logic. This mimics a scenario where only the attacker is present on the FPGA.

\textbf{Xilinx Ring Oscillators:}
Ring Oscillators are a malicious circuit with the sole purpose of aggressively consuming power. 
These are implemented as banks of combinational loops, resulting in rapid switching and power consumption as the circuit cannot settle on a single output value. 
Such a circuit can cause voltage disruptions in the power distribution network and can be used as a covert channel or to induce faults~\cite{gnad2017voltage, krautter2018fpgahammer, sugawara2019oscillator}.
%Classifying these circuits is important for detecting and removing these malicious circuits from multi-tenant and even single-tenant systems.
%These circuits do not pass Design Rule Checks (DRCs) and, in this simplistic implementation, cannot be deployed to the AWS platform. 
%As a result we consider an alternative power waster design based on the DSP components performing a pipe-lined fused multiple add operation synchronized from a clock source. For each platform we maximize the DSP utilization of the design respective of the device's resources. We deploy this application to both the PYNQ and AWS platforms. 

\textbf{Xilinx Arithmetic-Heavy: }
FPGAs are particularly well suited for high-intensity signal processing tasks with arrays of digital signal processors (DSPs). As an approximation of these structures, we implement arrays of DSPs performing a pipelined fused multiply-add operation. All DSPs operate in a single clock domain and compute upon data generated by a randomly-seeded, linear-feedback shift register.

%particularly useful for high precision floating point computations, like neural networks
%As a result we consider an alternative power waster design based on the DSP components performing a pipe-lined fused multiple add operation synchronized from a clock source. For each platform, we maximize the design's DSP utilization respective to the device's resources. We deploy this application to both the PYNQ and AWS platforms. 
\textbf{Xilinx Cryptographic Cores:}
We study ten different implementations of cryptographic computations consisting of two algorithms (AES, PRESENT) implemented on five different architectures (Custom HLS IP core and as software running on Orca, MicroBlaze, PicoRV, and ARM CortexM3 soft processors). %This adds ten unique applications. Orca-AES, Orca-PRESENT, MicroBlaze-AES, MicroBlaze-PRESENT, PicoRV-AES, PicoRV-PRESENT, Cortex-AES, Cortex-PRESENT, HLS-AES, and HLS-PRESENT.

\textbf{Intel NIOS-V/M:}
We implement a NIOS-V/M soft CPU core on the DE10-Nano to perform sensor characterization. We do not perform classification or other side-channel attacks. Instead, we provide a more thorough description of our novel phase-step techniques and describe the effects of this technique on $\theta$ and $\phi$ tuning. To tune $\phi$, we implement an AES-128 workload on the NIOS-V/M core. When enabled, the core will repeatedly perform decryption on a static plaintext.

\begin{comment}


%We study the channel capacity of our sensor on nine designs. We implement the applications within the same fabric, but logically isolated from our \name TDC. Seven of these applications are cryptographic cores, including four different architectures (custom IP core, ORCA, MicroBlaze, and PicoRV soft-processors) and two algorithms (AES and PRESENT). This yields the following applications: AES, ORCA-AES, ORCA-PRESENT, MB-AES, MB-PRESENT, PICO-AES, PICO-PRESENT. The remaining two applications are baseline, the absence of co-located computation, and power wasters, an array of combinatorial loops designed to consume power maliciously. The applications all operate at 25MHz.




One of these is a baseline, i.e., only the sensor is operating. Another is a power waster, which emulates a potential fault attack. We also look at many different implementations of cryptographic operations, including two different algorithms (AES and PRESENT) running on different architectures (custom IP core, Microblaze, ORCA, and PicoRV). We describe each of these applications in more detail.

\subsubsection{Baseline}
The primary goal of the baseline is to model the lack of another co-tenant.
The baseline only contains the \sensor and associated data collection logic.
This mimics a scenario where only the attacker is present on the FPGA.

\subsubsection{Power Waster}
A power waster is a malicious circuit with the sole purpose of aggressively consuming power.
This design implements banks of ring oscillators that can be turned on to consume significant amounts of power.
This causes voltage disruptions in the PDN and can be used as a covert channel or to induce faults.
Classifying these circuits is valuable for detecting these malicious computations in a multi-tenant hardware platform.

\subsubsection{Cryptographic Cores}
We study seven different implementations of cryptographic computations. These include two algorithms (AES, PRESENT) and four architectures (custom IP core and software running on  ORCA, MicroBlaze, and PicoRV soft processors). This provides seven more unique applications (HLS-AES, ORCA-AES, ORCA-PRESENT, MB-AES, MB-PRESENT, PICO-AES, PICO-PRESENT).
\end{comment}


%We study the channel capacity of our sensor on nine designs. We implement the applications within the same fabric, but logically isolated from our \name TDC. Seven of these applications are cryptographic cores, including four different architectures (custom IP core, ORCA, MicroBlaze, and PicoRV soft-processors) and two algorithms (AES and PRESENT). This yields the following applications: AES, ORCA-AES, ORCA-PRESENT, MB-AES, MB-PRESENT, PICO-AES, PICO-PRESENT. The remaining two applications are baseline, the absence of co-located computation, and power wasters, an array of combinatorial loops designed to consume power maliciously. The applications all operate at 25MHz.



% divide into base/PWs/HLS Cores/Soft-procs
%\subsubsection*{Baseline}
%The \name TDC, along with its support structure (MMCMs, memory interface, output storage, etc.) is placed independently on the FPGA fabric. The delay line is placed and routed uniformly as described in \ref{delay_line}. 

%\subsubsection*{Synchronous Power Waster}
%A power waster is a malicious circuit to aggressively consume power. Our design implements banks of ring oscillators that can be turned on to consume significant amounts of power and thus introduce disruptions in the power distribution network. In this implementation, the power waster is synchronous in that it is toggled based on an input clock. For the positive edge of the input clock, the power wasters are enabled, consuming relatively constant power for the duration of that edge. In the negative edge of the input clock, all the power wasters are disabled, consuming no power. The input clock is configured to be $25$MHz. The \name TDC is implemented as in {\it Baseline} and logically isolated from the synchronous power waster circuit.

%\subsubsection*{ORCA Soft-Processor}
%\subsubsection*{Soft-Processor Applications}
%The ORCA Soft-Processor is a RV32IM implementation of the RISC-V ISA~\cite{orca_vectorBlox}.
%It implements a four-stage pipeline without instruction or data caches on bare metal, i.e., without an operating system kernel or MMU.
%The software runs is a C implementation of AES-128~\cite{tiny_aes}.
%The core repeatedly runs this encryption operation and reuses its key, input plain text, and initialization vector.
%The processor operates at $25$MHz and is logically isolated from the sensor.

